\subsection{Nonsmooth contact and multibody time-stepping}
Classical rigid-contact solvers can be broadly divided into two lines: penalty/compliant contact and nonsmooth contact based on complementarity or differential variational inequalities. Jean's nonsmooth contact dynamics, the Stewart--Trinkle implicit time-stepping scheme, and the LCP/constraint-stabilization work of Anitescu--Potra and Anitescu--Hart collectively established the modern classical foundation of rigid contact based on velocity-level time stepping, unilateral constraints, and complementarity modeling \cite{jean1999nonsmooth,stewart1996implicit,anitescu1997formulating,anitescu2004constraint,anitescu2010iterative,tasora2011matrixfree,pfeiffer2012nonsmooth}. More recent tutorials and surveys have provided broader syntheses of contact constraints in Lagrangian systems, rigid-body frictional contact, and multibody contact mechanics \cite{contacttutorial2016,contactconstraints2016,flores2022contact,wang2021nonsmooth,peng2022sensitivity}.

Building on this foundation, work on specific solvers, contact patches, flexible multibody systems, quasistatic stepping, and numerical comparisons has further expanded the solution landscape of this line of research \cite{melanz2017comparison,catchingup2018,contactpatch2018,symplecticcontact2019,flexiblecontact2020,convexquasistatic2021,effectivemass2022,contactbasics2023,predictioncorrection2025,contactmultiscale2025}. On the implementation and application side, studies on flexible multibody strategies for large lap joints, rigid/compliant comparisons, wheel--rail conformal contact, automatic contact detection, and improved polygonal contact have advanced engineering usability from the perspectives of software realization, contact representation, and detection pipelines \cite{lapjoints2016,pazouki2017compliant,conformalwheel2020,autocontact2023,polygonalcontact2023,nencioni2025conformal,gradientcontact2025}. In robotics control and contact-rich manipulation, rigid contact and friction have also been incorporated explicitly into inverse dynamics and convex time-stepping formulations \cite{inversedynamics2016,convexquasistatic2021}.

It is important to emphasize that the LCP/MLCP line does not correspond exactly to ``all real contact,'' but rather to an idealized rigid-body nonsmooth contact model: bodies are treated as rigid, contact is represented by unilateral non-penetration constraints, and local elastic compression and contact-patch evolution are not resolved explicitly; instead, velocity or impulse updates satisfy complementarity conditions over discrete time steps \cite{jean1999nonsmooth,stewart1996implicit,anitescu1997formulating,contacttutorial2016,flores2022contact,wang2021nonsmooth}. As a result, this line most naturally describes constraint response, impulse transfer, and contact-state transitions in discrete time rather than the full local mechanics of contact in a continuum sense. Accordingly, LCP-based methods are most appropriate for rigid-body, multi-contact, impulse-dominated problems with frequent contact/separation switching. In near-conformal or large-area contact, local elastic compression and Hertz-type contact-patch evolution, lubrication- or adhesion-dominated contact, rate- and history-dependent friction, and cases requiring accurate contact timing, impact waveforms, or smooth local force histories, the results should be interpreted as equivalent constraint responses at the discrete-time level rather than as complete descriptions of real contact mechanics \cite{pazouki2017compliant,contactpatch2018,nencioni2025conformal}. In this sense, the literature has largely explained why complementarity-based rigid contact is meaningful, but it still lacks a systematic discussion of \emph{how local geometry should enter the contact model within this framework}.

\subsection{Implicit geometry, distance fields, and collision surrogates}
In geometry representation and collision detection, boundary-distance-based convex queries have long been shaped by GJK, Baraff's rigid-contact analysis, and dynamic BVH techniques, while cloth and deformable-body collision detection has further highlighted the importance of robust contact handling for dynamic quality \cite{gilbert1988fast,baraff1989analytical,baraff1994fast,larsson2006dbvh,bridson2002robust,teschner2005collision,wang2021collision}. On the other hand, distance fields have long served as implicit geometric surrogates for shape representation, sphere tracing, sparse voxel organization, and high-resolution narrow-band representations. ADF, sphere tracing, OpenVDB, and NanoVDB represent key developments in adaptive sampling, robust querying, and sparse data structures for distance fields \cite{frisken2000adf,hart1996sphere,museth2013vdb,museth2021nanovdb}.

For contact and collision processing based on SDFs, prior studies have investigated point-to-SDF continuous collision detection, locally optimized robust SDF collision handling, contact between beams/flexible bodies and discrete SDFs, SDF-DEM modeling for arbitrary granular shapes, SDF-enhanced CFD-DEM, real-time collision detection between general SDFs, and unified formulations bridging SDF contact and traditional discrete-element contact \cite{hapticsdf2017,macklin2020sdfcollision,sdfbeam2020,sdfdem2022,sdfcfddem2023,liu2024sdfsdf,lai2024unifying,lopez2024convexsdf}. More recent work has further discussed piecewise-continuous interpolation for SDFs, moving-SDF contact formulations, SDF-based conformal contact detection, dynamic SDF reconstruction for irregular deformable bodies, and the possibility of promoting SDFs to general proxies for robotic multi-contact \cite{adaptivesdf2023,movingframesdf2024,superellipsecontact2025,dynamicreconsdf2025,contactsdf2025}. These studies clearly demonstrate the value of SDFs as geometric surrogates, but they focus primarily on geometric queries, collision detection, compliant/proximal-style contact formulations, or generic contact proxy capability. In contrast, \emph{how local SDF geometry can improve contact modeling while preserving the main complementarity-based rigid hard-contact solution chain} remains insufficiently articulated. At the same time, interpolation, differencing, and nearest-point branch switching can amplify discontinuities in normals and curvature at the sub-voxel level, which means that the manner in which SDF geometry enters the dynamics model must itself be designed carefully.

\subsection{Local geometric representation and contact stabilization}
Computer graphics and geometric simulation have also long recognized the limitations of first-order contact approximation and have developed a variety of position-level, compliant-level, and patch-level stabilization strategies. Representative examples include Position Based Dynamics, XPBD, position-based solid simulation with improved contact accuracy, small-step explicit stabilization, and extended PBD variants for rigid bodies \cite{muller2007pbd,macklin2016xpbd,francu2017solids,macklin2019smallsteps,xpbdrigidbody2020}. Another class of work starts from asynchronous or variational contact and seeks to improve robustness and parallelism in complex contact through finer event handling or more consistent energy-based contact models \cite{harmon2009acm,vouga2011avcm,ainsley2012spacm}.

In optimization-based, proximal, and monolithic contact modeling, proximal-operator approaches, IPC, FrictionalMonolith, intersection-free rigid-body dynamics, contact-and-friction courses for graphics, monolithic frictional contact treatment for rigid/fluid systems, and recent unified handling of contact, friction, and shock propagation offer alternatives to the classical velocity-level LCP complementarity chain from the perspectives of unknowns, stabilization pathways, and coupling structure \cite{proximalcontact2017,li2020ipc,frictionalmonolith2021,contactgraphicscourse2021,intersectionfree2021,sphcontact2023,unifiedshock2023,chen2024primaldual}. In contrast to these position-level, optimization-level, or explicit proximal lines, the present work is concerned with a different question: how local geometric representation should enter the contact model without rewriting the underlying complementarity-based rigid-contact solution chain. Specifically, we do not treat the SDF as a universal contact proxy whose superiority must be established independently; rather, we use it as a carrier of local geometric information to construct patch-level local representations, generate curvature feedforward, and adopt different local modeling modes under preconfigured regimes corresponding to different contact states.
